\chapter{结果与讨论}

\section{示例结果}

表\ref{tab:result}使用三线表，表号和表题位于表的上方。表文采用小五号宋体；单位宜集中写入表头，避免在每个数据单元中重复。

\begin{table}[htbp]
  \centering
  \caption{示例验证结果}
  \label{tab:result}
  \ReportTableText
  \begin{tabular}{lccc}
    \toprule
    指标 & 目标值 & 观测值 & 判定 \\
    \midrule
    完整率/\% & $\geq 99.0$ & 99.6 & 通过 \\
    平均响应时间/ms & $\leq 200$ & 163 & 通过 \\
    重复试验次数 & $\geq 30$ & 30 & 通过 \\
    \bottomrule
  \end{tabular}
\end{table}

\section{不确定性与限制}

结果解释应区分观测事实、推断和评价结论。对于未达到预期的结果，应说明原因、影响范围和处置方式；不影响正文连贯性的明细数据可移至附录，但每个附录均应在正文中被提及，见附录 A。
